\documentclass[11pt, a4paper]{letter}
\usepackage[margin=2.5cm]{geometry}
\usepackage{hyperref}
\usepackage{parskip}

\signature{Saeyong Park}
\address{Saeyong Park\\
Avikus (HD Hyundai)\\
Republic of Korea\\
\texttt{saeyong.park@avikus.ai}}

\begin{document}

\begin{letter}{%
The Editors\\
\textit{Ocean Engineering}\\
Elsevier
}

\opening{Dear Editors,}

We submit for your consideration the manuscript entitled
\textbf{``WaveSplat: Grid-Free Spectral Wave Field Estimation from Sparse
Shipboard LiDAR''} for publication in \textit{Ocean Engineering}.

\textbf{Scope and motivation.}
Real-time wave field estimation from shipboard sensors is a critical
enabler for autonomous maritime operations, including proactive motion
compensation, route optimisation, and wave-preview feedforward control.
Existing methods rely on spatial grid discretisation combined with ensemble
Kalman filtering, which introduces interpolation artefacts and fixed spatial
resolution.
This submission addresses these limitations with a fundamentally different
representation.

\textbf{Principal contributions.}
WaveSplat represents the ocean wave field as a superposition of
spectral components — plane waves constrained by the linear dispersion
relation — optimised directly on sparse, irregularly distributed LiDAR
observations without any spatial grid.
Four novel contributions are presented:
\begin{enumerate}
  \item A \emph{grid-free spectral estimation framework} with persistent
    cross-frame tracking via dispersion-based phase prediction, enabling
    real-time operation with approximately 10 gradient steps per frame.
  \item A \emph{dispersion-constrained parameterisation} that reduces the
    free parameter count from $5N$ to $4N$ and acts as a physics-based
    regulariser against overfitting, quantified through ablation.
  \item An \emph{incidence-angle-weighted observation model} that accounts
    for the heteroscedastic noise structure of LiDAR measurements on ocean
    surfaces.
  \item An \emph{adaptive model complexity} mechanism that automatically
    selects the number of spectral components from $N{=}1$ for simple swells
    to $N{=}21$ for complex multi-component seas.
\end{enumerate}

\textbf{Key results.}
On a physics-based simulation dataset of 100 episodes spanning Beaufort 2--5
sea states (1--30 wave components), WaveSplat achieves:
\begin{itemize}
  \item \textbf{30.1\% lower normalised RMSE} than a hyperparameter-tuned
    Ensemble Transform Kalman Filter (ETKF) baseline
    (0.223 vs.\ 0.319, Wilcoxon $p{=}1.1\!\times\!10^{-15}$,
    Cohen's $d{=}1.16$);
  \item \textbf{approximately $10\times$ faster} inference
    (13.4 vs.\ 132.6 ms/frame, single CPU thread);
  \item the \textbf{only positive Skill Score} in a pre-registered
    six-metric evaluation panel against a zero-predictor reference,
    exposing both the ETKF (shape-incoherent) and sliding-window GPR
    (near-zero predictor) as failing this basic test.
\end{itemize}
The improvement is highly significant across 100 paired runs with five
random seeds, and is robust to both diffuse and specular LiDAR observation
models and to extreme observation sparsity ($<$8\% degradation at 20\%
observation density).

\textbf{Relevance to Ocean Engineering.}
This work directly addresses the operational challenge of real-time
wave field reconstruction for vessel-scale maritime applications.
The method requires no pre-trained models, no offline calibration beyond
water depth, and runs on standard CPU hardware at 10-Hz LiDAR frame rates.
Results include a companion application to wave-preview feedforward
motion control of hydrofoil vessels, demonstrating the practical value
of the spectral representation for downstream maritime control systems.

\textbf{Originality statement.}
This work is original and is not under consideration elsewhere.
The simulation code, all experiment scripts, and five representative
episodes are openly available at the project repository to ensure
end-to-end reproducibility of all reported tables and figures.

\textbf{Suggested reviewers.}
We suggest the following experts in ocean wave sensing and data
assimilation (the authors have no conflict of interest with any of them):
\begin{itemize}
  \item Prof.\ Tianhao Liu (wave sensing from ship-mounted sensors)
  \item Prof.\ Johannes Gemmrich (ocean surface wave observations)
  \item Dr.\ Nicolas Desmars (LiDAR-based wave field reconstruction)
\end{itemize}

We look forward to your consideration of this manuscript.

\closing{Sincerely,}

\end{letter}
\end{document}
